\begin{lemma}
\label{lemma-homotopy-category-abelian-split}
Let $\mathcal{A}$ be an abelian category. The following are equivalent.
\begin{enumerate}
\item The category $K(\mathcal{A})$ is abelian.
\item Every morphism of $K(\mathcal{A})$ satisfies condition (3) of
Lemma \ref{lemma-when-split}.
\item Every short exact sequence in $\mathcal{A}$ splits.
\end{enumerate}
If these conditions hold, then the functor
$$
H^* : K(\mathcal{A}) \longrightarrow K(\mathcal{A}), \qquad
H^*(X^\bullet)^n = H^n(X^\bullet), \quad d_{H^*(X)}^n = 0
$$
is an equivalence of categories.
\end{lemma}

\begin{proof}
By Proposition \ref{proposition-homotopy-category-triangulated},
$K(\mathcal{A})$ is pre-triangulated. Hence Lemma
\ref{lemma-when-split} and
Homology, Definition \ref{homology-definition-abelian-category}
show that (1) and (2) are equivalent. More explicitly, an abelian
category has kernels, so the lemma gives the normal form in (2).
Conversely, that normal form supplies kernels and cokernels, and the
canonical map from coimage to image becomes the identity on its middle
summand.

Assume (2), and let
$$
0 \longrightarrow A \xrightarrow{a} B \longrightarrow C
\longrightarrow 0
$$
be a short exact sequence in $\mathcal{A}$. Apply (2) to $a$, viewed as
a map between complexes concentrated in degree zero. The normal form gives
a morphism $b : B \to A$ in $K(\mathcal{A})$ such that $aba = a$.
There are no nonzero homotopies between maps of complexes concentrated in
degree zero, so this is the same equality in $\mathcal{A}$. Since $a$ is
a monomorphism, $ba = 1$, and the sequence splits. Thus (2) implies (3).

Assume (3). For a complex $X^\bullet$, put
$$
Z_X^n = \Ker(d_X^n), \qquad B_X^n = \Im(d_X^{n - 1}).
$$
Choose splittings of the two short exact sequences
$$
0 \longrightarrow Z_X^n \longrightarrow X^n \longrightarrow
B_X^{n + 1} \longrightarrow 0,
\qquad
0 \longrightarrow B_X^n \longrightarrow Z_X^n \longrightarrow
H^n(X^\bullet) \longrightarrow 0.
$$
They give decompositions
$$
X^n = B_X^n \oplus H^n(X^\bullet) \oplus B_X^{n + 1}
$$
under which $d_X^n(b, h, c) = (c, 0, 0)$. Consequently $X^\bullet$
is the direct sum of $H^*(X^\bullet)$ and the complex $C_X^\bullet$ with
$$
C_X^n = B_X^n \oplus B_X^{n + 1}, \qquad
d_{C_X}^n(b, c) = (c, 0).
$$
The maps $s^n : C_X^n \to C_X^{n - 1}$ given by
$s^n(b, c) = (0, b)$ form a contracting homotopy. We obtain maps
$$
i_X : H^*(X^\bullet) \longrightarrow X^\bullet, \qquad
p_X : X^\bullet \longrightarrow H^*(X^\bullet)
$$
whose composites are inverse in $K(\mathcal{A})$ and which induce the
identity on cohomology.

The maps on cohomology define $H^*$ on morphisms of $K(\mathcal{A})$ by
Homology, Lemma
\ref{homology-lemma-map-cohomology-homotopy-cochain}. For any map
$u : H^*(X^\bullet) \to H^*(Y^\bullet)$, the map
$i_Y u p_X$ induces $u$ on cohomology. Conversely, for every
$f : X^\bullet \to Y^\bullet$ in $K(\mathcal{A})$ we have
$$
f = i_Y p_Y f i_X p_X = i_Y H^*(f) p_X.
$$
Thus $u \mapsto i_Y u p_X$ is an inverse to the map on Hom sets induced
by $H^*$, so $H^*$ is fully faithful. It is essentially surjective since
every $X^\bullet$ is isomorphic to $H^*(X^\bullet)$ and
$H^*H^*(X^\bullet) = H^*(X^\bullet)$.

The full subcategory of complexes with zero differential is the category
$\text{Gr}(\mathcal{A})$ of
Homology, Definition \ref{homology-definition-graded}, which is abelian by
Homology, Lemma \ref{homology-lemma-graded}. The equivalence just proved
therefore implies that $K(\mathcal{A})$ is abelian. This proves (3)
implies (1) and finishes the proof.

The splittings used above are choices; no natural isomorphism between the
identity functor and $H^*$ is asserted.
\end{proof}
